\documentclass[11pt]{article}
\usepackage[margin=1in]{geometry}
\usepackage{amsmath,amssymb,amsthm}
\usepackage[T1]{fontenc}
\usepackage{lmodern}
\usepackage{microtype}
\usepackage[hidelinks,hypertexnames=false]{hyperref}
\newcommand{\Ec}{\mathbb E_c}
\newcommand{\Vc}{\operatorname{Var}_c}
\newcommand{\qh}{\widehat q}
\newcommand{\Zh}{\widehat Z}
\newtheorem{theorem}{Theorem}
\title{T25(iv): Conditional evidence variance under PW2}
\author{SMC-BVS v3 proof run}
\date{22 October 2026}
\begin{document}
\maketitle
\noindent\textbf{Scope.} This document proves only the conditional assertion
in T25(iv), including four-island pooling. It does not establish PW2 for
unrestricted inputs. All notation and hypotheses are those of
\texttt{STATEMENT.md}. The proof is not an adoption of the result.

\section{Setting and precise conclusion}
Fix the data and a successful frozen calibration value $c$. The finite
state space is $\Omega=\{S\subseteq[p]:|S|\le K\}$, with
$K=\min(p,\lfloor n/2\rfloor)$. The method's positive bridge is
$\gamma_\varphi=\gamma_0 e^{\varphi d}$, with
$\bar Z_\varphi=\sum_{S\in\Omega}\gamma_\varphi(S)$ and
$q_\varphi=\gamma_\varphi/\bar Z_\varphi$.
Here $\gamma_0(S)=\varpi(S)\prod_b r_b(S\cap C_b)$,
$r_b(U)=m(U)/m(\emptyset)$, and
$d(S)=\log[m(S)/m(\emptyset)]-\sum_b\log r_b(S\cap C_b)$.
Consequently $\gamma_1(S)=\varpi(S)m(S)/m(\emptyset)$.

Write $N=N_e$ for the number of particles in one island and fix
$0=\varphi_0<\cdots<\varphi_T=1$. Put
\[
 q_t=q_{\varphi_t},\qquad Z_t=m(\emptyset)\bar Z_{\varphi_t},\qquad
 G_t=e^{(\varphi_t-\varphi_{t-1})d},\qquad
 c_t=\frac{Z_t}{Z_{t-1}},\qquad g_t=\frac{G_t}{c_t}.
\]
Thus $Z_T=Z$, $q_{t-1}(G_t)=c_t$, and, for every function $f$,
\begin{equation}\label{eq:bridge}
 q_{t-1}(g_t f)=q_t(f).
\end{equation}
The empirical measure $\qh_t$ is taken after stage-$t$ resampling and
mutation. The particles start iid from $q_0$, with the exact normaliser
$\Zh_0=Z_0$. The evidence update and its normalised form are
\begin{equation}\label{eq:path}
 \Zh_t=\Zh_{t-1}\qh_{t-1}(G_t),\qquad
 W_t:=\frac{\Zh_t}{Z_t}
       =W_{t-1}\qh_{t-1}(g_t),\qquad W_0=1.
\end{equation}
Each stage resamples multinomially and mutates with the method's
$q_t$-invariant kernel $K_t$. All settings are fixed before final
randomness (CAL), and all operations use exact arithmetic and fresh ideal
streams (EX-A, EX-S). The four final islands use independent streams.
No conditioning on final diagnostics is performed.

\begin{theorem}[T25(iv), conditional form]\label{thm:main}
For each $t=1,\ldots,T$, let $a_t,\nu_t$ be finite nonnegative constants
fixed given $c$. Suppose that the actual island satisfies
\begin{equation}\tag{PW2}\label{eq:pw2}
 \Ec\!\left[W_{t-1}^2
       \{(\qh_{t-1}(g_t))^2-1\}\right]
 \le \frac{a_t}{N}\Ec W_{t-1}^2+\nu_t.
\end{equation}
Then
\begin{equation}\label{eq:bound}
 \Vc(W_T)\le B_T:=
 \prod_{t=1}^T\left(1+\frac{a_t}{N}\right)-1
 +\sum_{s=1}^T\nu_s
       \prod_{t=s+1}^T\left(1+\frac{a_t}{N}\right).
\end{equation}
If each of the four final islands satisfies PW2 with these same constants,
then for $W_{\rm pool}=\frac14\sum_{r=1}^4 W_{r,T}=\Zh/Z$,
\begin{equation}\label{eq:poolbound}
 \Vc(W_{\rm pool})\le\frac{B_T}{4}.
\end{equation}
Empty products are one. Here $N=N_e$ throughout.
\end{theorem}

\section{Finiteness and the first moment}
Every $d(S)$ is finite. Write $d_- =\min_\Omega d$ and
$d_+=\max_\Omega d$. Since the grid is increasing,
\[
 e^{\varphi_t d_-}\le \frac{\Zh_t}{Z_0}
                       \le e^{\varphi_t d_+}.
\]
Also $Z_t>0$ is deterministic given $c$. Therefore $W_t$ has finite
moments of every order for each fixed $c$ and $t$. In particular all
expectations in PW2 and the following calculations are well defined.

Wave-1 T12(i.a), transferred to conditional calibration in T12(iii),
gives $\Ec[\Zh_t\qh_t(f)]=Z_tq_t(f)$. For completeness, its needed
first-moment argument follows directly from the prescribed operations.
Let $\mathcal F_{t-1}$ include the island's history through stage
$t-1$. Multinomial resampling and mutation give
\begin{equation}\label{eq:condmean}
 \mathbb E_c[\qh_t(f)\mid\mathcal F_{t-1}]
 =\frac{\qh_{t-1}(G_t K_t f)}{\qh_{t-1}(G_t)}.
\end{equation}
Indeed each child first chooses a parent with weights proportional to
$G_t$ and then moves with $K_t$; averaging its conditional expectation
over the $N$ children proves the identity.

We prove simultaneously for all functions $f:\Omega\to\mathbb R$ that
\begin{equation}\label{eq:weighted}
 \Ec[W_t\qh_t(f)]=q_t(f).
\end{equation}
At $t=0$ this follows from iid $q_0$ starts and $W_0=1$. For the
induction step, $W_t$ in \eqref{eq:path} is
$\mathcal F_{t-1}$-measurable. Multiplying \eqref{eq:condmean} by it,
taking expectation and using the induction hypothesis with the
deterministic function $g_tK_tf$ gives
\begin{align*}
 \Ec[W_t\qh_t(f)]
 &=\Ec[W_{t-1}\qh_{t-1}(g_tK_tf)]\\
 &=q_{t-1}(g_tK_tf)=q_t(K_tf)=q_t(f).
\end{align*}
The last two equalities use \eqref{eq:bridge} and invariance. With
$f\equiv1$ this yields
\begin{equation}\label{eq:mean}
 \Ec W_t=1,\qquad \Vc(W_t)=\Ec W_t^2-1.
\end{equation}
This is an unconditional first moment under the frozen-run law; it does
not assert $\mathbb E_c[W_t\mid\mathcal F_{t-1}]=W_{t-1}$.

\section{Second-moment recursion and its solution}
Set $M_t=\Ec W_t^2$ and $\alpha_t=1+a_t/N\ge1$.
Squaring the pathwise update \eqref{eq:path} and subtracting
$W_{t-1}^2$ gives the exact identity
\begin{equation}\label{eq:increment}
 M_t-M_{t-1}
 =\Ec\!\left[W_{t-1}^2
       \{(\qh_{t-1}(g_t))^2-1\}\right].
\end{equation}
Applying PW2 to this identity yields
\begin{equation}\label{eq:rec}
 M_t\le\alpha_t M_{t-1}+\nu_t,\qquad M_0=1.
\end{equation}
By \eqref{eq:increment}, PW2 is equivalent to the recursion
$M_t\le(1+a_t/N)M_{t-1}+\nu_t$. The theorem therefore adds no control of
the evidence beyond PW2 itself; all of its substance lies in establishing
PW2 with useful constants, which is not done here.
Define
\[
 U_0=1,\qquad
 U_t=\prod_{j=1}^t\alpha_j
       +\sum_{s=1}^t\nu_s\prod_{j=s+1}^t\alpha_j.
\]
Separating the last term of the sum shows
$U_t=\alpha_t U_{t-1}+\nu_t$. Induction now proves $M_t\le U_t$:
the base case is equality, and if $M_{t-1}\le U_{t-1}$ then
\[
 M_t\le\alpha_tM_{t-1}+\nu_t
       \le\alpha_tU_{t-1}+\nu_t=U_t,
\]
where the second inequality uses $\alpha_t\ge0$.
By \eqref{eq:mean}, $\Vc(W_T)=M_T-1\le U_T-1=B_T$.
This proves \eqref{eq:bound}, with the slack at stage $s$ multiplied
only by the factors from later stages $t>s$.

\section{Four independent islands}
Under CAL and EX-S, conditional on $c$, each island is a measurable
function of the fixed settings and its own fresh ideal stream. Thus the
four evidence ratios are independent and identically distributed given
$c$, as also established in wave-1 T12(iv.c). Each has mean one by
\eqref{eq:mean}, and each has variance at most $B_T$. Hence
\[
 \Vc\!\left(\frac14\sum_{r=1}^4W_{r,T}\right)
 =\frac1{16}\sum_{r=1}^4\Vc(W_{r,T})
 \le\frac{4B_T}{16}=\frac{B_T}{4}.
\]
The factor $1/4$ applies to both terms in \eqref{eq:bound}.
It is not obtained by replacing $N_e$ with $4N_e$ inside that formula.
Terminal rounds and local search leave the evidence unchanged, so this
is the bound for the reported pooled evidence. Since the argument holds
for each successful frozen calibration, it holds almost surely given
$\mathcal C$ on calibration success. This completes the proof.
\hfill$\square$

\paragraph{Dependency record.}
The wave-1 source is \texttt{COLLECTED\_PROOFS\_v2.tex}, chapter
``T12 (T8): frozen-run evidence and recorded Rao--Blackwell identities'',
parts (i.a), (iii), (iv.c). These are used as proved, not adopted results;
the needed first-moment induction is also displayed above. No unproved
v3 conjecture is invoked. In particular, PW2 remains a hypothesis.
\end{document}
